\documentclass[11pt]{amsart}
\usepackage{amsmath,amssymb,amsthm,mathtools}
\usepackage[margin=1in]{geometry}
\usepackage{enumerate}
\usepackage{hyperref}

\newtheorem{theorem}{Theorem}[section]
\newtheorem{proposition}[theorem]{Proposition}
\newtheorem{lemma}[theorem]{Lemma}
\newtheorem{corollary}[theorem]{Corollary}
\theoremstyle{definition}
\newtheorem{definition}[theorem]{Definition}
\newtheorem{problem}[theorem]{Problem}
\newtheorem{principle}[theorem]{Meta-principle}
\theoremstyle{remark}
\newtheorem{remark}[theorem]{Remark}
\newtheorem{observation}[theorem]{Observation}

\newcommand{\RH}{\mathrm{RH}}
\newcommand{\bbR}{\mathbb{R}}
\newcommand{\bbC}{\mathbb{C}}
\newcommand{\bbN}{\mathbb{N}}
\newcommand{\bbQ}{\mathbb{Q}}
\newcommand{\bbZ}{\mathbb{Z}}
\newcommand{\dm}{dm_\infty}
\newcommand{\dmfull}{dm_{\mathrm{full}}}
\newcommand{\dmon}{dm_{\mathrm{full,on}}}
\newcommand{\dnu}{d\nu}
\newcommand{\Wl}{W_\lambda}
\newcommand{\Tl}{T_\lambda}
\newcommand{\negind}{\mathrm{neg.ind}}
\DeclareMathOperator{\Real}{Re}
\DeclareMathOperator{\Imag}{Im}

\title[Exhaustion of 1-D Propagation]{The exhaustion of one-dimensional propagation:\\
a structural no-go for the Riemann Hypothesis,\\
and two directions that survive it}

\author{David Alejandro Trejo Pizzo}
\thanks{Independent Researcher. \texttt{dtrejopizzo@gmail.com}}


\begin{document}
\nocite{bgConnes99,bgDeninger}

\begin{abstract}
A research programme spanning thirty-five phases has reduced the Riemann
Hypothesis (RH), under the Connes--Consani--Moscovici (CCM) framework, to
the vanishing of a single non-negative measure $d\nu=\dmfull-\dmon$, and
has catalogued seven structural walls (MW-1 through MW-7) that block every
attempted proof. We make two contributions. \emph{First (the no-go):} we
isolate the common logical form of all seven walls. Every live or recently
closed route presents RH as the \emph{boundary point} of a one-parameter
family whose every other member is an unconditional theorem (the
De~Bruijn--Newman time $t$, the spectral-triple approximation parameter
$x$, the place-set $S$, the Abel scale $\lambda$). We argue that the walls
are precisely the statement that \emph{no one-dimensional propagation
principle reaches that boundary}, and we verify this by exhibiting, for
each classical 1D mechanism --- limit/continuity, analyticity, parabolic
backward uniqueness, monotonicity/integrality, ergodic recurrence --- the
exact point at which it either is excluded by a wall or coincides with RH
itself. As a sharpening internal to the framework we prove that the Abel
kernel $\Wl$ is \emph{strictly positive at every real point}, whence the
universal quantifier in the trace criterion ``$\Tl=0$ for all $\lambda$''
is redundant: a single scale suffices, and the $\lambda$-family collapses
to a point. \emph{Second (the survivors):} we show the no-go leaves
exactly two logical forms unexcluded --- \emph{two-dimensional propagation}
(coupling the heat time with the arithmetic filtration via a maximum or
monotonicity principle on a quadrant) and \emph{index amplification} (a
non-grouplike tensor operation on the Pontryagin space $(K,Q)$ realising
the Deligne mechanism: integer invariant, multiplicative defect,
sublinear uniform bound). We give each a precise functional specification,
identify what would falsify it, and state the open problems. We do not
solve them.
\end{abstract}

\maketitle
\tableofcontents
\newpage

\section{Introduction}

\subsection{What this paper is, and is not}
This is not a paper that proves the Riemann Hypothesis, nor one that
claims partial progress toward a proof. It is a paper about the
\emph{shape} of the obstruction. After a long programme of attempts---each
documented, each refuted or reduced to a known wall---we are in the
unusual position of possessing a near-complete map of why RH resists the
methods currently available. A complete map of failure is itself a
mathematical object, and it has a use: it constrains, by elimination, the
logical form that any successful argument must take. The purpose of this
paper is to extract that constraint rigorously and to deduce from it where
the search should go next.

We work throughout inside the CCM framework \cite{CCM24}, in the form
developed in \cite{P39,P40}. We recall the minimum needed in
\S\ref{sec:framework}. The two halves of the paper are independent in
spirit: \S\ref{sec:metapattern}--\S\ref{sec:nogo} are a structural no-go
argument (with one new unconditional theorem, \S\ref{sec:wlambda});
\S\ref{sec:formA}--\S\ref{sec:formC} are a forward research agenda with
falsifiable specifications but no theorems claiming to reach RH.

\subsection{The one new theorem}
The only genuinely new unconditional result proved here is
Theorem~\ref{thm:wlambda-strict}: the Abel kernel $\Wl$ of the CCM trace
criterion is strictly positive at every real point (not merely almost
everywhere), and consequently the trace criterion of \cite{P39} holds at a
\emph{single} scale $\lambda$ rather than requiring all $\lambda>0$. This
is modest in difficulty---a three-term-recurrence argument---but it is
exactly the kind of statement the structural analysis predicts should be
available: it removes a spurious degree of freedom from the map without
touching the hard core.

\section{The framework, recalled}\label{sec:framework}

We fix notation following \cite{P39,P40}.

\begin{itemize}
\item The archimedean spectral measure is
$\dm(s)=(2\pi)^{-2}\,|\Gamma(1/4+is/2)|^2\,ds$ on $\bbR$, a probability
measure (after normalisation) with density decaying like
$|s|^{-1/2}e^{-\pi|s|/2}$.
\item $\{\mathcal{P}_n\}_{n\geq0}$ are the orthonormal polynomials for
$\dm$, with $\mathcal{P}_0\equiv1$ and Jacobi data $b_n=0$ (parity),
$a_n^\infty=\tfrac12\sqrt{(2n+1)(2n+2)}$. They obey the three-term
recurrence
\begin{equation}\label{eq:ttr}
s\,\mathcal{P}_n(s)=a_n^\infty\,\mathcal{P}_{n+1}(s)+a_{n-1}^\infty\,\mathcal{P}_{n-1}(s),
\qquad a_k^\infty>0.
\end{equation}
\item $\dmfull=|\zeta(1/2+is)|^2\,\dm$ and $\dmon=|\zeta_{\mathrm{on}}(1/2+is)|^2\,\dm$,
where $\zeta_{\mathrm{on}}$ is the Hadamard product of $\zeta$ restricted to
its on-line zeros.
\item The central object is the difference measure
$\dnu=\dmfull-\dmon$. By \cite[]{P40} it is \emph{unconditionally a
non-negative measure}: $\dnu=(R-1)\,|\zeta_{\mathrm{on}}|^2\,\dm$ with
$R=|\zeta/\zeta_{\mathrm{on}}|^2\geq1$, each off-critical zero quadruple
contributing a factor $\geq1$.
\item The Abel kernel is, for $N=N(\lambda):=\#\{n:\gamma_n\leq\lambda\}$,
\begin{equation}\label{eq:wlambda}
\Wl(s)=\sum_{n=1}^{N} n\,|\mathcal{P}_n(s)|^2+(a_N^\infty)^2\,|\mathcal{P}_{N+1}(s)|^2\;\geq0.
\end{equation}
\item The trace is $\Tl=\int_\bbR \Wl\,\dnu$. The CCM trace criterion
\cite[Thm]{P39} states
\begin{equation}\label{eq:p39}
\RH\iff \Tl=0\ \text{for all }\lambda>0.
\end{equation}
\item \cite[]{P40} lists eight equivalent reformulations of RH as
functionals of $\dnu$ ($\dnu=0$; the Cauchy transform $\mathcal{G}\equiv0$;
$\Tl=0$; the Jensen defect $\delta=0$; the variational distance
$\tilde\delta_\lambda^2=0$; $\mathrm{KL}=0$; $d_{\mathrm{CCM}}=0$).
\end{itemize}

\begin{definition}[The Hadamard barrier, MW-7]\label{def:hadamard}
No property of $\zeta$ that is verifiable without knowledge of the
positions of its zeros implies $\dnu=0$.
\end{definition}

The Hadamard barrier is the framework-internal face of the master
obstruction. Its companions across the programme are: MW-1 (Weil
positivity $\equiv$ RH), MW-2 (the arithmetic propagation wall: the Euler
product converges only for $\Real(s)>1$ and its information does not
propagate to the critical strip; zeros are global objects), MW-3 (no
Hasse--Minkowski in infinite dimension), MW-4 (every unconditional tool
yields a positivity bound of the wrong sign), MW-5 (the arithmetic-site
factorisation requires the unfinished Connes--Consani programme), MW-6
(uniform spectral-gap control $\equiv$ RH). The thesis of
\S\ref{sec:metapattern}--\S\ref{sec:nogo} is that these seven are seven
faces of one fact.

\section{The meta-pattern: RH as a boundary point}\label{sec:metapattern}

\subsection{The recurring shape}
Across the programme, every route to RH---whether eventually closed or
still live---instantiates one template. There is a one-parameter family of
mathematical objects; for all parameter values except one boundary value
the relevant statement is an unconditional theorem; and the boundary value
is exactly RH. We tabulate the instances, with their sources.

\begin{center}
\renewcommand{\arraystretch}{1.3}
\begin{tabular}{llll}
\textbf{Family} & \textbf{Parameter} & \textbf{Unconditional control} & \textbf{Boundary $\iff$ RH}\\
\hline
Self-approximation $\mathrm{SR}_d$ & $d$ & proved for all $d\neq0$ & $d=0$\\
DBN flow $\Tl(t)$ & $t$ & $\Tl(t)=0$ for $t\geq\Lambda$ & $t=0$\\
Spectral triples $H_x$ & $x$ & real spectrum, all finite $x$ & $x=\infty$\\
Semilocal index $\mu_S$ & $S$ & positivity for small $S$ & $S=$ all places\\
Pontryagin $\kappa_\lambda(N)$ & $N$ & controllable, all finite $N$ & $N=\infty$\\
Euler product & $\sigma$ & converges for $\sigma>1$ & $\sigma=1/2$\\
\end{tabular}
\end{center}

The sources are, respectively: the Bagchi/Voronin self-approximation
family (proved for all $d\neq0$ by Baker, Nakamura, Pa\'nkowski,
Garunk\v{s}tis; the closing of $d=0$ is RH \cite[]{P40}); the
De~Bruijn--Newman flow with $\Lambda$ the DBN constant
\cite[Docs~70--71]{P40}; the CCM finite-Euler spectral triples
\cite{CCM25} analysed in \cite[]{P40}; the semilocal Weil
positivity of Connes--Consani \cite{CC21} analysed in \cite[]{P40};
the Pontryagin negative-index packet \cite[]{P40}; and the classical
Euler product itself.

\begin{principle}[The boundary-point pattern]\label{prin:boundary}\sloppy
In every analysed approach, RH is the limit/boundary member of a
one-parameter family whose interior is unconditionally controlled. The
structural walls MW-1--MW-7 are, jointly, the assertion that no
one-dimensional propagation principle carries the interior control to the
boundary.
\end{principle}

We do not claim Meta-principle~\ref{prin:boundary} as a theorem; it is an
empirical regularity of the corpus, stated to be tested. \S\ref{sec:nogo}
gives it teeth by checking the propagation mechanisms one by one.

\section{Sharpening: the Abel family collapses to a point}\label{sec:wlambda}

Before the no-go we record the one place where the boundary-point pattern
degenerates favourably: the Abel scale $\lambda$ is not a genuine
parameter at all.

\begin{lemma}[No shared zeros of consecutive orthogonal polynomials]\label{lem:noshared}
For every $n\geq0$, the polynomials $\mathcal{P}_n$ and $\mathcal{P}_{n+1}$
have no common zero in $\bbC$.
\end{lemma}
\begin{proof}
Suppose $\mathcal{P}_{n+1}(s_0)=\mathcal{P}_n(s_0)=0$. Equation
\eqref{eq:ttr} at index $n$ gives
$0=a_n^\infty\mathcal{P}_{n+1}(s_0)+a_{n-1}^\infty\mathcal{P}_{n-1}(s_0)
=a_{n-1}^\infty\mathcal{P}_{n-1}(s_0)$, so $\mathcal{P}_{n-1}(s_0)=0$ since
$a_{n-1}^\infty>0$. Descending induction yields
$\mathcal{P}_0(s_0)=0$, contradicting $\mathcal{P}_0\equiv1$.
\end{proof}

\begin{theorem}[Strict pointwise positivity of $\Wl$]\label{thm:wlambda-strict}
For every $\lambda$ with $N(\lambda)\geq1$, $\Wl(s)>0$ for all $s\in\bbR$.
\end{theorem}
\begin{proof}
By \eqref{eq:wlambda}, $\Wl$ is a sum of non-negative terms with strictly
positive coefficients. If $\Wl(s_0)=0$ then every term vanishes; in
particular the term $n=N$ (present because $N\geq1$) forces
$\mathcal{P}_N(s_0)=0$, and the last term forces
$\mathcal{P}_{N+1}(s_0)=0$. This contradicts Lemma~\ref{lem:noshared}.
\end{proof}

\begin{remark}
This strengthens \cite[Cor.~7.5]{P40}, which established
$\Wl>0$ off a finite set (the nodes of $\mathcal{P}_{N+1}$) and neutralised
that set using the absolute continuity of $\dnu$. The exceptional set is in
fact empty, and the conclusion needs no hypothesis on $\dnu$. For the
degenerate range $N(\lambda)=0$ one has $\Wl=(a_0^\infty)^2|\mathcal{P}_1|^2$,
vanishing only at $s=0$, and the original argument applies.
\end{remark}

\begin{theorem}[A single scale suffices]\label{thm:single-scale}
Fix any $\lambda_0>0$ with $N(\lambda_0)\geq1$. Then $\RH\iff\Tl[\lambda_0]=0$.
The universal quantifier in \eqref{eq:p39} is redundant.
\end{theorem}
\begin{proof}
($\Rightarrow$) Under RH, $\dnu=0$, so $\Tl[\lambda_0]=0$. ($\Leftarrow$)
Write $\Tl[\lambda_0]=\int \Wl[\lambda_0]\,(R-1)\,|\zeta_{\mathrm{on}}|^2\,\dm$
with $R-1\geq0$ (\cite[]{P40}) and $\Wl[\lambda_0]>0$ everywhere
(Theorem~\ref{thm:wlambda-strict}). The integrand is non-negative; its
integral vanishing forces $(R-1)|\zeta_{\mathrm{on}}|^2=0$ $\dm$-a.e.; since
$|\zeta_{\mathrm{on}}|^2>0$ off a null set, $R\equiv1$ a.e., i.e.\ $\dnu=0$,
i.e.\ RH (\cite[]{P40}, condition 2).
\end{proof}

The Abel family is thus not a propagation problem at all: its boundary is
already glued to its interior. This is the exception that proves the rule.
For every \emph{other} family in the table, the boundary is genuinely
detached, and \S\ref{sec:nogo} explains why no 1D principle can reattach it.

\section{The no-go: one-dimensional propagation is exhausted}\label{sec:nogo}

We now run through the classical mechanisms by which an unconditional
interior statement could be propagated to a boundary value, and show each
is either excluded by a named wall or is logically equivalent to RH. The
list is intended to be exhaustive over the 1D toolkit; we invite the reader
to supply a mechanism not reducible to one of these entries.

\subsection{Limit / continuity}
If the family $f(\theta)$ were continuous at the boundary and identically
zero on the interior, the boundary value would vanish by continuity. For
$H_x\to H_\infty$ this is exactly the route: Hurwitz's theorem applied to
the regularised determinants would transfer reality of zeros from finite
$x$ to $x=\infty$ \cite[Prop.~5.3]{P40}. But the quantitative core
of that limit is $\lim\mu_\lambda\geq0$, the semilocal Weil positivity in
the limit---this is MW-1 wearing the mask of MW-6. Continuity to the
boundary \emph{is} the prize, not a tool for it.

\subsection{Analyticity}\label{ssec:analyticity}
This deserves a clean statement because it explains a recurring
frustration.

\begin{proposition}[No faithful analytic detector]\label{prop:analytic}
Suppose $F:[0,\infty)\to\bbR$ is real-analytic and satisfies $F(t)=0$ if
and only if the zeros of $\Xi_t$ are all real. Then RH holds.
\end{proposition}
\begin{proof}
The DBN flow has $\Xi_t$ with only real zeros for all $t\geq\Lambda$
(definition of the DBN constant $\Lambda$, with $\Lambda\geq0$ known
unconditionally and $\Lambda\leq0$ being RH). Hence $F\equiv0$ on
$[\Lambda,\infty)$. A real-analytic function vanishing on a set with a
limit point vanishes identically, so $F\equiv0$, whence $F(0)=0$, i.e.\ the
zeros of $\Xi_0=\Xi$ are all real: RH.
\end{proof}

Proposition~\ref{prop:analytic} shows that constructing a real-analytic
faithful detector of zero-reality is \emph{already} a proof of RH. This is
why the natural detector $\Tl(t)$ is non-analytic precisely at $t=\Lambda$:
the on/off split that defines $\dmon$ jumps there. Analyticity is not an
available tool; it is the trophy. Any candidate ``analytic monotone
quantity'' must be checked against this proposition---if it is faithful and
analytic, the work is already done, which means it cannot be both faithful
and constructible without circularity.

\subsection{Parabolic backward uniqueness}
The DBN flow is a (backward) heat flow, and backward uniqueness for
parabolic equations would propagate $u(\Lambda)=0\Rightarrow u\equiv0$.
This works for quantities \emph{linear} in $\Xi_t$. The reality-of-zeros
detector is nonlinear (it sees the imaginary parts of zeros, not a linear
functional of $\Xi_t$). This is the ``linear $\neq$ nonlinear'' gap
identified as the terminus of Path~III in \cite{P40}; backward uniqueness
is available exactly where it is useless.

\subsection{Monotonicity / integrality}
The Pontryagin negative index $\kappa(Q)=\negind(H_C)$ is an integer
(\cite{P35}, with $\kappa=2m$, localised in $\mathcal{K}_{\mathrm{off}}$ by
\cite[Thm~8.1]{P40}) and upper semicontinuous. If a local principle
forbade its jumps along a deformation to the on-line configuration,
integrality would force $\kappa=0$. But the jumps occur exactly at zero
collisions, and nothing local prohibits them---generic heat flow
\emph{creates} complex pairs on crossing $\Lambda$ for other entire
functions. Integrality constrains the value but not its constancy.

\subsection{Ergodic recurrence}
Bagchi's criterion makes RH equivalent to strong recurrence of the
vertical flow $\tau\mapsto\zeta(\cdot+i\tau)$ on the Bohr torus
$\Omega=\prod_p\mathbb{T}_p$. The Kronecker flow is uniquely ergodic, so the
arithmetic point $\omega_0$ is uniformly recurrent unconditionally; the
obstruction is not recurrence but the continuity of the Euler observable
$\Phi(\omega)=\prod_p(1-\omega(p)p^{-s})^{-1}$ at the arithmetic point,
which holds iff $\sum_p p^{-s}$ converges in the strip, iff RH
\cite[Prop.~7.2]{P40}. The ergodic route is MW-2 in dynamical
language, with no residue.

\subsection{Conclusion of the no-go}
\begin{observation}\label{obs:nogo}
Each classical 1D propagation mechanism, applied to the boundary-point
families of \S\ref{sec:metapattern}, either coincides with RH
(continuity \S5.1, analyticity \S5.2) or is excluded by a named wall
(backward uniqueness by nonlinearity, integrality by collision jumps,
ergodicity by MW-2). The 1D toolkit is exhausted.
\end{observation}
Consequently, an argument that reaches RH cannot be a one-dimensional
propagation principle. By elimination, the candidate logical forms that
remain are: \emph{(A)} two-dimensional propagation, where information flows
across a region from its sides to a corner without passing continuously
through any single boundary point; and \emph{(B)} amplification, where a
defect is multiplied until it collides with an unconditional sublinear
bound. The rest of the paper specifies these.

\section{Form A: two-dimensional propagation}\label{sec:formA}

\subsection{Why a second dimension is not excluded}
Every wall in \S\ref{sec:nogo} concerns a \emph{1D} boundary. MW-2 says the
$x$-axis (arithmetic filtration) alone does not reach the corner; the DBN
analysis says the $t$-axis (heat time) alone does not. \emph{No result in
the corpus addresses their coupling.} Two-dimensional propagation
principles---the maximum principle, Almgren-type monotonicity, logarithmic
convexity along diagonals---move information from the \emph{edges} of a
region into its interior and corners, and do not require continuity through
any single point. This is the gap.

\subsection{The candidate object}
Define the joint trace
\begin{multline*}
\mathcal{T}(t,x)=\text{(CCM trace of the finite-Euler approximant indexed by }x\text{,}\\
\text{evolved by the DBN heat flow to time }t).
\end{multline*}
Its edges and corner are controlled by results already in hand:
\begin{itemize}
\item Edge $\{t\geq\Lambda\}$: $\mathcal{T}=0$ unconditionally (DBN, \cite[Docs~70--71]{P40}).
\item Edge $\{x\text{ finite}\}$: real spectrum unconditionally (CF rigidity of the CCM triples, \cite[]{P40}, which passed the circularity filter I2a at finite level).
\item Corner $(t,x)=(0,\infty)$: RH.
\item $\partial_t\mathcal{T}$ is computed: a quadratic form in the Jacobi data $c_k$, with exact cancellation under RH \cite[]{P40}.
\item $\partial_x\mathcal{T}$ is computed in discrete form: primes enter one at a time through $\Tl=A_\lambda^{\mathrm{off}}-\sum_p(\log p/\sqrt p)B_\lambda(\log p)$ \cite[]{P40}.
\end{itemize}

\begin{problem}[The coupling inequality]\label{prob:coupling}
Does $\mathcal{T}(t,x)$ satisfy a differential inequality coupling
$\partial_t$ and $\partial_x$---parabolic, monotone, or logarithmically
convex along diagonals of the quadrant---under which positivity propagates
from the two edges to the corner $(0,\infty)$, even though it fails along
each axis separately?
\end{problem}

\begin{remark}[What would falsify Form A]
If $\partial_t\mathcal{T}$ and $\partial_x\mathcal{T}$ have no fixed sign
relation on the quadrant---equivalently, if the mixed second variation
$\partial_t\partial_x\mathcal{T}$ changes sign with no diagonal along which
a monotone quantity is built---then no maximum/monotonicity principle
applies and Form~A is dead. The first test is finite and honest: compute
$\partial_t\mathcal{T}$ for the two-prime approximant and compare its sign
to $\partial_x\mathcal{T}$.
\end{remark}

\section{Form B: index amplification (the Deligne mechanism)}\label{sec:formB}

\subsection{The skeleton of the one proof of RH that exists}
RH is a theorem over function fields \cite{Del74}. Stripped to its logical
skeleton, Deligne's mechanism is not a positivity argument:
\begin{enumerate}[\rm(1)]
\item an \emph{integer} invariant (the weight of a Frobenius eigenvalue);
\item an \emph{amplification}: a tensor/symmetric power multiplies the
defect, $\kappa\mapsto k\kappa$;
\item a \emph{uniform sublinear bound} ($k\kappa\leq C(k)$ with $C(k)/k\to0$,
from rationality of the family's $L$-function: the poles do not move);
\item conclusion: $k\kappa\leq C(k)$ for all $k$ forces $\kappa=0$.
\end{enumerate}
None of (1)--(4) is a CAP-type positivity. The programme audited the
\emph{cohomological formalism} of Weil--Deligne (it is MW-5); it did not
audit the \emph{amplification trick} as an abstract skeleton. The skeleton
is what we transplant.

\subsection{What the programme already supplies}
Ingredient (1) is in hand: $\kappa(Q)=\negind(H_C)$ is an integer,
$\kappa=2m$, localised in $\mathcal{K}_{\mathrm{off}}$
(\cite{P35}; \cite[]{P40}). The missing ingredient is the
amplification.

\begin{problem}[The amplifying functor]\label{prob:amplify}
Does there exist an operation $\boxtimes$ on the Pontryagin space $(K,Q)$,
or on the category of CCM spectral triples, such that
\begin{enumerate}[\rm(a)]
\item $\kappa(\psi_1\boxtimes\psi_2)\geq\kappa(\psi_1)+\kappa(\psi_2)$
(superadditivity of the negative index), and
\item there is a bound $\kappa(\psi^{\boxtimes k})\leq f(k)$ with
$f(k)/k\to0$, provable by a finiteness condition of type C5
(\cite[]{P40})?
\end{enumerate}
If both hold, then $2mk\leq f(k)$ for all $k$ forces $m=0$, i.e.\ RH.
\end{problem}

\subsection{Why it has not been tried, and where to look}
For $\zeta$---a $\mathrm{GL}(1)$ object---the \emph{grouplike} tensor theory
is trivial: $\mathrm{sym}^k$ of the trivial character is trivial, which is
why the Langlands route to Ramanujan-type bounds works for modular forms
but says nothing about $\zeta$. The operation $\boxtimes$ must therefore be
\emph{non-grouplike}: a candidate is the $\lambda$-ring / Frobenius-lift
structure on the Connes--Consani scaling site, or a Rankin--Selberg-type
convolution realised directly on $(K,Q)$ rather than on automorphic data.

\begin{remark}[What would falsify Form B]
If $\kappa$ is exactly additive under every available $\boxtimes$ (not
strictly superadditive) and the bound $f(k)$ is genuinely linear, there is
no contradiction to extract and Form~B is dead. The decisive sub-question
is whether \emph{any} natural product on $(K,Q)$ makes the negative index
strictly superadditive.
\end{remark}

\begin{remark}[Resonance with the one non-CAP mechanism found]
The single mechanism in the entire corpus that is provably not a CAP
positivity is Bagchi's amplification \cite[]{P40}: one off-critical
zero, replicated by strong recurrence to a positive-density family,
contradicts the unconditional zero-density bound $N(\sigma,T)=o(T)$.
Form~B is that same skeleton with an algebraic amplifier in place of the
ergodic one. The ergodic version failed only because its premise (strong
recurrence) is RH-equivalent (\S\ref{sec:nogo}); the skeleton itself is
sound. This is the strongest single piece of evidence that the
amplification \emph{form} is the right place to look.
\end{remark}

\section{Form C: the boundary point as a quantitative singularity}\label{sec:formC}

A cheaper, lower-risk line. The self-approximation family $\mathrm{SR}_d$ is
proved unconditionally for all $d\neq0$, with $d=0\iff$RH
\cite[]{P40}. The effective question has not been asked.

\begin{problem}[Effective degeneration of $\mathrm{SR}_d$]\label{prob:effective}
As $d\to0$, how do the constants in the proof of $\mathrm{SR}_d$ degrade?
If the degradation is $e^{-c/|d|}$, this is quantified MW-2 and the line
closes cleanly. If it is polynomial in $1/|d|$, there is a genuine tension
with the Rouch\'e argument linking recurrence to zero density, deserving
a dedicated analysis.
\end{problem}

Form~C produces knowledge either way: a closed line (with MW-2 quantified)
or a tension to exploit. It is the natural first move because it is finite,
falsifiable, and independent of the harder Forms A and B.

\section{Honest assessment}\label{sec:assessment}

We claim no proof and no partial proof of RH. We claim three things.
First, a structural no-go (Observation~\ref{obs:nogo}): the
one-dimensional propagation toolkit is exhausted against the
boundary-point families that organise the subject, and we have checked the
mechanisms individually. Second, one new unconditional theorem internal to
the framework (Theorem~\ref{thm:single-scale}): the trace criterion holds at
a single scale, the $\lambda$-family being spurious. Third, a forward
agenda (Problems~\ref{prob:coupling},~\ref{prob:amplify},~\ref{prob:effective})
that is precisely the set of logical forms the no-go does not exclude, each
with a stated falsifier.

If Forms A and B also collapse onto known walls, that collapse will itself
be informative: it would show the obstruction is deeper than the
1D/2D/amplification trichotomy can express, and would be the first
evidence that RH is unreachable by \emph{any} of the argument-forms
currently known to mathematics. That, too, would be a result. The honest
position is that we have narrowed the search for the missing idea from
``anything'' to ``a non-grouplike amplification of the Pontryagin index,
or a coupled monotonicity on the heat-time/filtration quadrant,'' and that
this narrowing---not any claim of progress toward the proof itself---is the
contribution.

\begin{thebibliography}{99}

\bibitem{CCM24} A.~Connes, C.~Consani, H.~Moscovici,
\emph{Spectral realisation of zeta zeros and the Riemann--Roch programme}
(2024). Foundational construction underlying this paper.

\bibitem{CCM25} A.~Connes, C.~Consani, H.~Moscovici,
\emph{Zeta spectral triples and finite Euler approximation},
arXiv:2511.22755 (2025); A.~Connes, arXiv:2602.04022 (2026).

\bibitem{CC21} A.~Connes, C.~Consani,
\emph{Weil positivity and trace formula for the archimedean place},
Selecta Math. \textbf{27} (2021), arXiv:2006.13771.

\bibitem{Del74} P.~Deligne, \emph{La conjecture de Weil. I},
Publ. Math. IH\'ES \textbf{43} (1974), 273--307.

\bibitem{P35} D.~A.~Trejo Pizzo, \emph{Pontryagin spectral index of the Weil form and the Connes scaling operator: a bridge theorem for the Riemann Hypothesis}, preprint, 2026.

\bibitem{P39} D.~A.~Trejo Pizzo, \emph{A trace and measure criterion for the Riemann Hypothesis from the zeta spectral triple}, preprint, 2026.

\bibitem{P40} D.~A.~Trejo Pizzo, \emph{The localized Weil form and its Kre\u\i n--Pontryagin index: localization, obstruction, and rigidity}, preprint, 2026.

\bibitem{bgConnes99} A.~Connes, \emph{Trace formula in noncommutative geometry and the zeros of the Riemann zeta function}, Selecta Math.\ (N.S.) \textbf{5} (1999), 29--106.
\bibitem{bgDeninger} C.~Deninger, \emph{Some analogies between number theory and dynamical systems on foliated spaces}, Doc.\ Math.\ ICM \textbf{I} (1998), 163--186.
\end{thebibliography}

\end{document}
